% LINA · POP, procedimento operacional padrão (entrada, meio, finalização).
% Compile com XeLaTeX. Só comandos de lina.cls: nada de \vspace, \fontsize, \color, \hspace ou \newpage aqui.
\documentclass[pop]{lina}

\referencia{POP-04}
\titulocurto{POP-04 · v. 2.1 · vigência 03/2026}
\siglarodape{LINA · POP-04}

\begin{document}

% ---------- Entrada ----------
\rotulo{Procedimento operacional padrão}
\titulo{Nome do procedimento, começando por um substantivo de ação}
\begin{metadados}[0]{YYY}
  \rotulometa{Responsável}\newline Coordenação técnica &
  \rotulometa{Aplica-se a}\newline edifício · objeto &
  \rotulometa{Duração}\newline \mono{≈ 4 h}
\end{metadados}

\secaofio{Objetivo}
Uma ou duas frases sobre o resultado que o procedimento garante, não sobre as tarefas.

\section{Material}
Equipamento · acessório · ficha de campo · mídia de cópia

\secaotabela{Procedimento}
\begin{passos}
  \item Cada passo começa com um verbo no imperativo e descreve uma única ação.
  \item Quando houver um critério de conferência, escreva-o no mesmo passo.
  \item Valores e parâmetros remetem a tabelas numeradas, não ao texto corrido.
  \item Nomes de arquivo e caminhos vão em mono: \mono{AAAA-MM-DD\_local\_E01}.
\end{passos}

\begin{atencao}
  \rotuloatencao{Atenção.} Uma caixa por página, no máximo, para riscos ao equipamento ou às pessoas.
\end{atencao}

\paginademo

% ---------- Meio ----------
\secaotabela{Procedimento \continuacao{(continuação)}}
\begin{passos}
  \item A numeração continua entre páginas; nunca reinicie a contagem.
  \item Quando um passo depender de condição, escreva a condição primeiro: “Se a prévia falhar, repita o passo 04.”
  \item Agrupe passos longos em etapas com subtítulo, mantendo a numeração corrida.
  \item O último passo é sempre o de registro: o que salvar, onde e com que nome.
\end{passos}

\titulotabela{Parâmetros por escala do objeto.}
\begin{tabela}{YW{142}W{130}}
  Escala & Resolução & Distância \\
  objeto & \numero{— mm} & \numero{— m} \\
  edifício & \numero{— mm} & \numero{— m} \\
  cidade & \numero{— mm} & \numero{— m} \\
\end{tabela}
\nota{Valores de referência; ajuste conforme o manual do equipamento.}

\secaoafastadatabela{Critérios de aceitação}
\begin{criterios}
  \item Cada estação compartilha ao menos três alvos com a anterior.
  \item O erro de registro não excede o valor da tabela 1.
  \item Todos os arquivos estão no servidor com o nome padrão.
\end{criterios}

\paginademo

% ---------- Finalização ----------
\section{Registro de execução}
Preencha uma linha por execução e arquive esta folha com os dados.
\begin{tabela}[12.7]{W{112}YW{152}W{90}}
  Data & Local / estações & Operador & Conferido \\
  \registro{5}
\end{tabela}

\secaoafastadatabela{Histórico de revisões}
\begin{tabela}[6.4]{W{82}W{122}Y}
  Versão & Data & Alteração \\
  1.0 & \numero{05/2024} & Emissão inicial. \\
  2.0 & \numero{08/2025} & Inclusão da tabela de parâmetros por escala. \\
  2.1 & \numero{03/2026} & Ajuste do padrão de nome de arquivo. \\
\end{tabela}

\secaoafastada{Documentos relacionados}
\numero{POP-03} Calibração do escâner · \numero{POP-07} Registro de nuvens de pontos

\vfill
\assinaturas{Elaborado por}{Revisado por}{Aprovado por}

\end{document}
